\begin{abstract}
We argue that grounding in dynamic environments requires a runtime architecture for observation and belief update, not just larger models or offline knowledge. LLMs are static compressors of syntax, physical plausibility, and historical facts, but agentic intelligence requires navigating situated states that cannot be learned from text. This grounding gap is structural: frozen parameters contain no guaranteed information about the current state. Bridging it requires an Agent Harness---a runtime architecture supplying observation, belief update, and execution. We formalize this framework via Bayesian inference, show the gap is bounded by observation fidelity rather than model scale, derive three design principles, and validate them empirically on a code-agent benchmark. Our analysis reframes the scaling debate from model capability to system architecture, showing that environmental adaptation comes from observation and action channel design, not parameter tuning.
\end{abstract}
% 中文翻译：
% 我们认为，动态环境中的落地需要一种用于观测和信念更新的运行时架构，而非仅仅依赖更大的模型或离线知识。LLM是语法、物理合理性和历史事实的静态压缩器，但真正的行动智能需要驾驭无法从文本中学习的情境化状态。这一落地鸿沟是结构性的：冻结参数不包含关于当前状态的任何保证信息。弥合它需要智能体框架——一种提供观测、信念更新和执行的运行时架构。我们通过贝叶斯推理形式化该框架，证明鸿沟受观测保真度而非模型规模约束，推导出三条设计原则，并在代码智能体基准上进行了实证验证。我们的分析将规模扩展的讨论从模型能力转向系统架构，表明环境适应来自观测与行动通道设计，而非参数调优。